\section{引言：最后一讲想回答什么}

这是 CS224N 的最后一讲。Manning 教授没有再去讲一条新的模型管线，而是把课程里最重要的线索重新串起来：NLP 的主线是什么、语言学在今天还重要吗、神经模型和符号系统到底是什么关系、以及我们应该如何看待 LLM 带来的机会与风险。

\begin{importantbox}{本讲的四个中心问题}
\begin{enumerate}
    \item CS224N 里那些最重要的想法，最后到底收束成了什么？
    \item NLP 和语言学的关系，是被 Transformer 时代削弱了，还是重新变得重要了？
    \item 语言的“意义”应该怎样理解，神经模型真的在做语义吗？
    \item 当模型变得更强之后，我们应该担心什么，应该继续研究什么？
\end{enumerate}
\end{importantbox}

\begin{figure}[H]
\centering
\includegraphics[width=0.95\textwidth]{slides-images/slide-000.jpg}
\caption{课程封面：CS224N 最后一讲，主题是 NLP、语言学与哲学\protect\footnotemark}
\end{figure}
\footnotetext{Slides 第1页。}

整门课可以粗略地归纳成一条连续的演进路线：
\begin{itemize}
    \item 从词向量和分布式语义开始；
    \item 到神经网络、序列模型、RNN/LSTM；
    \item 再到 Transformer、预训练和大模型；
    \item 然后进一步进入基准评测、推理、可解释性、多语言能力和安全问题。
\end{itemize}

\begin{knowledgebox}{课程里最稳固的几条经验}
\begin{itemize}
    \item dense representation 和上下文语义是现代 NLP 的基础。
    \item residual connection 让训练深层模型变得可靠。
    \item language modeling 看似简单，但确实是极强的通用预训练目标。
    \item benchmark 不是终点，generalization、reasoning、meaning 和 safety 仍然没有解决。
\end{itemize}
\end{knowledgebox}

\subsection{本章小结}

最后一讲不是“补一门新课”，而是把前面的技术路线提升成三个更大的问题：语言如何表示、语言如何理解、以及语言技术如何进入真实世界。

%% ========== 课程回顾 ==========

